\section{Introduction}
\label{sec:intro}

\subsection{What the program does}

\texttt{rubiks-cfop} does two things:

\begin{enumerate}[leftmargin=1.4em]
  \item it \textbf{generates} scrambles for the $3\times3\times3$ cube in the form
        the World Cube Association uses, and
  \item it \textbf{solves} a cube with \textbf{CFOP}: \emph{Cross}, \emph{F2L}
        (the first two layers, four corner--edge pairs), \emph{OLL} (orientation of
        the last layer) and \emph{PLL} (permutation of the last layer).
\end{enumerate}

The program is two layers thick.  A \emph{headless engine} does all of the work and
returns structured data: it never prints, never reads the environment and has no
third-party runtime dependencies.  A thin \emph{command line interface} turns
arguments into engine calls and data into text.  A graphical interface can be
added later without touching the engine, which is the subject of
Appendix~\ref{app:gui}.

\subsection{Conventions}

\begin{description}[leftmargin=0em,style=nextline]
  \item[Cube representation] Kociemba's facelet notation: 54 facelets ordered
  \texttt{U1..U9}, \texttt{R1..R9}, \texttt{F1..F9}, \texttt{D1..D9},
  \texttt{L1..L9}, \texttt{B1..B9}.  The solved cube is the string
  \texttt{UUUUUUUUURRRRRRRRRFFFFFFFFFDDDDDDDDDLLLLLLLLLBBBBBBBBB}.
  \item[Moves] WCA notation (Regulation 12a): the face turns
  \texttt{U D L R F B}, the wide turns \texttt{Uw Dw Lw Rw Fw Bw}, the rotations
  \texttt{x y z}, and --- as a documented extension --- the slice turns
  \texttt{M E S}.  Every move accepts the suffixes \texttt{2} and \texttt{'}.
  \item[Colors] the standard scheme of Regulation 4d1, white on top and green in
  front: \texttt{U} white, \texttt{R} red, \texttt{F} green, \texttt{D} yellow,
  \texttt{L} orange, \texttt{B} blue.
  \item[Cross] always solved into the \textbf{bottom layer}.  For a white cross the
  upper layer is yellow, and the cube is reoriented so that this is what you see.
\end{description}

\subsection{WCA compliance}

The scramble generator follows Regulation 4b3, which asks an official scramble
program to ``produce a random state from all states that require at least 2 moves
to solve (equal probability for each state)''.  For the $3\times3\times3$ cube that
means \emph{random-state} scrambles: a uniformly random cube state is drawn and then
solved, and the inverse of that solution is the scramble.  Random-move sequences ---
fast, filtered, and what many tools mean by ``a scramble'' --- are also available,
but they are not WCA-legal for this puzzle (Regulation 4b3e permits them only for
$5\times5\times5$ and larger).  Scrambling orientation follows Regulation 4d1.

\subsection{Where the algorithms come from}

The OLL (57 cases) and PLL (21 cases) and F2L (41 cases) algorithms in the JSON
tables were taken from the published \emph{Logiqx cubing-algs} collection, whose
case identifiers match the community numbering.  They are not trusted on the
engineer's word: every entry is \emph{verified by execution}, and the recognizer's
patterns are \emph{derived from the algorithms themselves}, so the table and the
recognition can never disagree.  Section~\ref{sec:tables} describes the mechanism.

\subsection{How this manual is organized}

Section~\ref{sec:install} installs the program, Section~\ref{sec:quickstart} shows
it working, Section~\ref{sec:cli} is the complete command line reference and
Section~\ref{sec:output} explains the output line by line.
Section~\ref{sec:notation} is a notation primer, Section~\ref{sec:solver} explains
how each stage is solved, Section~\ref{sec:tables} documents the JSON tables and
Section~\ref{sec:library} the Python API.  Section~\ref{sec:trouble} covers
troubleshooting and the known limitations, and Appendix~\ref{app:gui} plans the
graphical interface.
